\section{The heights of the period maps}\label{sec:heights}

Fix the weights, nodes and shift of \S\ref{sec:theta-chain}. The same
$h_u$ must bound the height of $B_u$ in \S\ref{sec:period-slopes} and
satisfy \eqref{eq:acct-height}. We compute the contribution of the CM
period to this bound, then treat the auxiliary prime and the base
lattice in \eqref{eq:acct-prime} and \eqref{eq:acct-base}.

\subsection{The CM period and the signatures}

Put $L=x^\varepsilon$, $a=xb$, and let $D_E$ denote the root
discriminant of $E$. For the differential of
\S\ref{sec:inert-kernel}, the doubled slope of each CM coordinate line is
\[
 \ell=\log\Omega=\gamma+\log(\pi/v)-\lambda_1,\qquad v=D_E/2.
\]
This includes the finite contribution $-\log(D_E/\sqrt d)$.
On the windows of Proposition~\ref{prop:field-window}, $\ell<0$ for
large $x$ and $|\ell|=o(L)$. Indeed $\sqrt d\le D_E\le4nqd$ by the
discriminant formula in the towers through $K$ and through
$\mathbb Q(\zeta_{4nqd})$. In the proof of that proposition
$|\lambda_1|\ll J\log U$, while $\log d\gg U^{\varepsilon/5}\log^2 U$
and $\log(4nqd)=o(L)$ on $[U,U^J]$ for fixed $J$.

At lift $j$ the prescribed signature is
$(p_j,q_j)=(k+j-1,m-j)$
\citep[\S15.14]{adriazola-grh}. Write $w_{j\rho}=w_j(\rho)$. Then
\begin{equation}\label{eq:height-signature}
 \sum_{j,\rho}(p_j-q_j-k)w_{j\rho}
      =\frac{f h_0m(m^2-1)}6=:\mathcal S .
\end{equation}
Indeed $p_j-q_j-k=2j-m-1$ has sum zero over $1\le j\le m$.
In \eqref{eq:chain-nodes} the term $b_1(\rho)$ drops out, and
$\sum_j(2j-m-1)(j-\tfrac12)=m(m^2-1)/6$. In particular
$\mathcal S\ell\le0$ on these windows.

\subsection{The height and the target slope}

Write $\mathbb E_\sigma$ for the average over all embeddings of the
coefficient field $F$. All operator norms below use the frames of the
factorization $B_u=P_uK_{m,u}\cdots K_{1,u}$. On $R_\sigma$ take a
reference form and an orthogonal decomposition into base types $z$.
In the true form let the squared norm on type $z$ be the reference
squared norm times $\kappa_\sigma\prod_j\lambda_{jz\sigma}$, with
positive factors and the same orthogonal decomposition. Put
\[
 c=\mathbb E_\sigma\log\kappa_\sigma,\qquad
 b_{\rm fr}=\mathbb E_\sigma\min_z\sum_j\log\lambda_{jz\sigma},
\]
where the minimum uses nonzero types. Submultiplicativity and the
inverse factors on $R_\sigma^\vee$ give, for $B_u\ne0$,
\begin{equation}\label{eq:height-product}
 \operatorname{ht}(B_u)\le
 \mathbb E_\sigma\log\|P_{u,\sigma}\|_{\rm ref}^2+
 \sum_j\mathbb E_\sigma\log\|K_{j,u,\sigma}\|^2-b_{\rm fr}-c .
\end{equation}
Indeed the true squared dual norm of a functional is the sum of its
reference squared norms on the types divided by the respective factors.

Let $\phi(y)=y\log y-y$ for $y>0$ and $\phi(0)=0$.
For a word $u$, let $\nu_{j\rho i}$ be the nonnegative row degrees
of the kernel polynomial and $b_{j\rho}$ the logarithmic hook factors
of its frame. Write $a_\rho^\circ=a_t$ at a top $(0,0,t)$ and zero
elsewhere. Suppressing $u$ in the subscripts, put
\begin{align*}
 n_{j\rho}&=\sum_i\nu_{j\rho i},&
 \Delta_u^n&=\sum_{j,\rho}(n_{j\rho}-a_\rho^\circ-k w_{j\rho}),\\
 \Delta_u^\phi&=\sum_{j,\rho}
       \left(\sum_i\phi(\nu_{j\rho i})-b_{j\rho}
                                      -\phi(a_\rho^\circ)\right).
\end{align*}

\begin{hypothesis}[\openinput]\label{hyp:height-norms}
For the prescribed kernels and the pairing on the lattices of
\S\ref{sec:inert-gain}, the frames above realize the following
bounds with nonnegative allowances whenever $B_u\ne0$.
\begin{align*}
 \mathbb E_\sigma\log\|K_{j,u,\sigma}\|^2
 &\le \sum_\rho\left(\sum_i\phi(\nu_{j\rho i})-n_{j\rho}\ell\right)
           +e_{K,j,u},\\
 \mathbb E_\sigma\log\|P_{u,\sigma}\|_{\rm ref}^2
 &\le -A\log T_u+e_{P,u},\qquad
 b_{\rm fr}\ge\sum_{j,\rho}b_{j\rho}-e_{{\rm fr},u}.
\end{align*}
For the common target $C=R^\vee$ in the same metric one has
\[
 \mu(C)-c-\mathcal L\le
       \ell\sum_{j,\rho}(p_j-q_j)w_{j\rho}+e_R,
\]
with the finite correction $\mathcal L$ of \S\ref{sec:inert-gain}.
The row degrees, hook factors and allowances are specified for all
words, with
\[
 e_R+e_{P,u}+e_{{\rm fr},u}+\sum_j e_{K,j,u}
                +\Delta_u^\phi+|\ell|\,(\Delta_u^n-\mathcal S)_+
       \ll_\varepsilon L\mathcal B
\]
uniformly on the windows.
\end{hypothesis}

The kernel comparison uses the normalized theta derivatives in the CM
coordinate lines. The target comparison requires the signed covolume
bound for the base types and their passage to $R$. The upper bound for
$\mu_{\max}(R)$ used in the cutoff of \S\ref{sec:period-slopes} controls
a different slope \citep[Hypotheses~15.10--15.12, 15.16, 15.18
and Theorem~15.70]{adriazola-grh}.

\begin{theorem}\label{thm:height-account}
Under Hypothesis~\ref{hyp:height-norms} there are $h_u$ with
$h_u+A\log T_u$ independent of $u$ and $\operatorname{ht}(B_u)\le h_u$
whenever $B_u\ne0$. They satisfy \eqref{eq:acct-height} with
$e_4\ll_\varepsilon L\mathcal B$.
\end{theorem}

\begin{proof}
Put $E_u^{\infty}=e_{P,u}+e_{{\rm fr},u}+\sum_j e_{K,j,u}$ and
\[
 t_u=-A\log T_u-c+E_u^{\infty}+
      \sum_{j,\rho}\left(\sum_i\phi(\nu_{j\rho i})-b_{j\rho}
                                                   -n_{j\rho}\ell\right).
\]
By \eqref{eq:height-product}, $\operatorname{ht}(B_u)\le t_u$ for
each nonzero $B_u$. Add the bound for $\mu(C)-c-\mathcal L$.
The top terms give $m\sum_t(\phi(a_t)-a_t\ell)=H_\Omega$, and
\eqref{eq:height-signature} gives exactly
\[
 \mu(C)+t_u-\mathcal L+A\log T_u-H_\Omega
 \le e_R+E_u^\infty+\Delta_u^\phi+(\mathcal S-\Delta_u^n)\ell .
\]
As $\ell<0$, the last term is at most $|\ell|(\Delta_u^n-\mathcal S)_+$.
Take $h_*=\max_u(t_u+A\log T_u)$ and $h_u=h_*-A\log T_u$.
These dominate the $t_u$ and give the same maximum on the left.
The positive part of that maximum is an admissible $e_4$.
\end{proof}

With this choice the spread required in \eqref{eq:period-error} is
\[
 U=\max_u\bigl(\mu_{\max}(N_u)+A\log T_u\bigr)
                      -\langle\mu(N_u)+A\log T_u\rangle.
\]
The common $h_*$ cancels. The integral inclusions and the frame
estimates of \S\S\ref{sec:period-slopes}--\ref{sec:bessel-periods}
are still required on these lattices.

\subsection{The auxiliary prime}

Use the full column dimension $n_u$ and the formal degree budget $D^*$
of \S\ref{sec:bessel-periods}. Put $r=k+m-1$, $g=\log(g_0g_1)$ for
the two auxiliary level primes, and
\[
 d_u^*=\operatorname{mass}(k)I_\Pi
       \prod_{\rho\in\Phi}\dim(\lambda^{\rm col}_{u,\rho}),\qquad
 X_u=(D^*/d_u^*)^{1/(fk)} .
\]
Here $\operatorname{mass}(k)$ is the mass of the unimodular hermitian
genus, $I_\Pi$ is the column level index and $\dim$ is Weyl's formula
on the shifted highest weight.

\begin{hypothesis}[\openinput]\label{hyp:height-dimensions}
The row and column parameters of the chosen words satisfy the spacing,
interlacing and size conditions (P1)--(P4), the mass formula and the
target dimension comparison of
\citep[Proposition~15.25]{adriazola-grh} with the nodes
\eqref{eq:chain-nodes} and this $D^*$. For their prescribed primes
$T_u$ the full columns satisfy
\[
 \frac{D^*}{C_0\Gamma_0^2}\le n_u\le d_u^*(T_u^k-1)^f,\qquad
 \log\Gamma_0\le fk,\qquad \log T_u\ll_\varepsilon mL ,
\]
with $C_0\ge1$ fixed and $\Gamma_0=\Gamma_0(x)\ge1$ as in the column count of
\S\ref{sec:bessel-periods}.
\end{hypothesis}

\begin{proposition}\label{prop:height-prime}
Under Hypothesis~\ref{hyp:height-dimensions}, \eqref{eq:acct-prime}
holds with $e_6\ll_\varepsilon L(1+m^2/k)$.
\end{proposition}

\begin{proof}
The cited Weyl and mass comparison gives
\[
 \log X_u\ge\sum_j\log(\bar w_j v/\pi)+mg
                              -O_\varepsilon\bigl(L(1+m^2/k)\bigr).
\]
In that comparison the factors $\prod_{s=k}^{r-1}s!$ in the consecutive
masses cancel the denominators in the ratio of Weyl dimensions.
The size conditions $m\le x$, $k\le x^{O(1)}$ and
$\pi\le v\le e^L$ hold here by \eqref{eq:arithmetic-ranks} and the
discriminant bounds above. Since $n_u\le d_u^*T_u^{fk}$, its lower
bound gives
\[
 \log T_u\ge\log X_u-\frac{\log C_0+2\log\Gamma_0}{fk}.
\]
The last loss is at most $2+(\log C_0)/(fk)$.
\end{proof}

\subsection{The base lattice}

Write $\lambda_\rho=(\lambda_{\rho1},\ldots,\lambda_{\rho k})$ for
the dominant level-zero weight of a word and
$\nu_\rho=\lambda_\rho-\lambda_{\rho k}(1,\ldots,1)$.
For a weight $\lambda$, write $|\lambda|=\sum_i\lambda_i$, including
its signs.

\begin{hypothesis}[\openinput]\label{hyp:height-base}
The actual $M_u$ of \S\ref{sec:inert-gain} is the orthogonal sum over
lattice classes $c'$ of $\bigotimes_{\rho\in\Phi}
\mathbb S_{\lambda_\rho}Q_{c',\rho}$ with the integral Weyl lattices
and the Fock norms normalized at the highest weight vector. The
$Q_{c',\rho}$ are metrized base changes of rank-$k$ hermitian
$\mathcal O_E$-lattices $L_{c'}\subseteq L_{c'}^\#$ whose hermitian
duals satisfy $[L_{c'}^\#:L_{c'}]=T_u^{2f}$. The $\lambda_\rho$
have the signed weight bounds of \S\ref{sec:theta-chain}.
\end{hypothesis}

\begin{proposition}\label{prop:height-base}
Under Hypotheses~\ref{hyp:height-dimensions} and~\ref{hyp:height-base},
\eqref{eq:acct-base} holds with
$e_7\ll_\varepsilon N_T aL(1+m/k)$.
\end{proposition}

\begin{proof}
Duality negates the slope. The index contributes
$2\log(T_u^{2f})/([E:\mathbb Q]k)$ to the slope difference, hence
$\mu(Q_{c',\rho})=-\log T_u/k$. The determinant character of
$\mathbb S_\lambda$ in rank $k$ is
$\det^{|\lambda|\dim(\mathbb S_\lambda)/k}$, as one sees on scalar
matrices. Taking determinants, then tensor products and the
orthogonal sum of equal-slope terms, gives
\[
 \mu(M_u)=-\frac{\log T_u}{k}\sum_\rho|\lambda_\rho|
                       +\sum_\rho\mu_{\rm std}(\nu_\rho).
\]
Here $\mu_{\rm std}(\nu)$ is the slope of the Weyl module of weight
$\nu$ on the standard rank-$k$ lattice with its normalized Fock norm.
Its lower bound is
\[
 \mu_{\rm std}(\nu)\ge
       -|\nu|\log k-\tfrac12k(k-1)\log(\nu_1+k)
\]
by the hook determinant and the integral Gram bound
\citep[Lemma~15.29 and Proposition~15.30]{adriazola-grh}.

At a top, the first entry is $a_t$ and the last is $-b_t$. All other
entries there and all entries elsewhere are $O(b')$, since the common
bound in \S\ref{sec:theta-chain} is at most the first node at any top,
which is $O(b')$. Thus both $\sum_\rho(|\lambda_\rho|)_+$ and
$\sum_\rho|\nu_\rho|$ are $O(N_Ta+fkb')=O(N_Ta)$, since
$b'=a/Z$ and $Z$ dominates every fixed power of $x$. Likewise
$fk^2\log(a+k)\ll N_TaL$. Insert these estimates and
$\log T_u\ll_\varepsilon mL$ in the determinant formula.
\end{proof}

Since $A\le2N_Ta$ and $\mathcal B\ge N_Ta(1+m^2/k)$, the three
estimates give
\[
 e_4+A e_6+e_7\ll_\varepsilon x^\varepsilon\mathcal B
\]
under Hypotheses~\ref{hyp:height-norms}, \ref{hyp:height-dimensions}
and~\ref{hyp:height-base}. The finite correction $E_u$ in
\eqref{eq:acct-local} remains to be estimated for the same words.
